\section{Experimental Design}
\label{sec:experiments}

\subsection{Evidence protocol}

Experiments were executed through repository scripts and preserved as immutable
run directories. Report tables were extracted from machine-readable JSON and CSV
records rather than copied from console output. Formal A/B/C evidence is kept
separate from earlier technical smokes and from the edge-width sub-ablation. This
separation is necessary because the candidate sources and experimental questions
differ.

The principal asset is an AK-47 mesh with 26,133 triangles, 21,548 UV coordinates,
928 UV islands, and 4,999 welded-topology seam pairs. The accepted texture size is
2048 by 2048 RGB for final export. Formal comparison runs use 512-pixel working
textures for controlled evaluation and share the target mesh, cameras, metrics,
palette, candidate count, and base seeds.

\subsection{Formal A/B/C comparison}

The formal experiment contains four paired candidates with seeds 20260716 to
20260719:

\begin{description}
  \item[\routeA] a generic pattern-first square source with periodic repair;
  \item[\routeB] a whole-weapon composition created in continuous weapon space,
  baked through OBJ-derived geometry maps, and corrected at UV edges;
  \item[\routeC] the corresponding \routeB\ result plus one diagnosis and two local
  correction attempts, subject to the 0.01 improvement gate.
\end{description}

\routeA\ and \routeB\ intentionally use different route-specific generation logic.
The comparison asks whether the weapon-aware representation produces better
mapped behavior, not whether identical pixels change after a single processing
operation. Within each route, the target, seed schedule, resolution, cameras, and
evaluation remain fixed.

For paired differences, the report gives the mean, median, standard deviation, win
rate, and bootstrap 95\% confidence interval. With only four pairs, these intervals
describe the observed candidate set and should not be read as population-level
proof across arbitrary themes or weapons.

\subsection{Seam-correction and selection-policy sub-ablations}

The seam-correction table compares each formal \routeB\ texture immediately before
and after the locked correction. This isolates the effect of the seam operation
without changing source content.

A separate edge-width sweep uses one unchanged earlier technical composition and
tests widths 0, 2, 4, 8, 12, 16, and 24 pixels. The hard conditions are
$\assetseam\leq0.01$ and multi-view retention of at least 95\%. The same candidate
pool is ranked by:

\begin{enumerate}
  \item feasibility, Pareto dominance, and deterministic objective order; and
  \item an equal 50/50 normalized weighted score.
\end{enumerate}

This sweep supports claims about operating-point and policy behavior. It is not
used to claim that the earlier artwork is visually superior to the final
generated candidates.

\subsection{Three-checkpoint live case}

The interaction case begins with a dragon theme and the AK-47 adapter. The Agent
produces four directions; the user selects \emph{Treasured Relic}. Four landscape
artwork variations are generated and mapped at low cost. One failed source
attempt, containing large empty regions and isolated clusters, is rejected and
regenerated locally rather than restarting the complete run. The user selects
\texttt{artwork\_04} after viewing the source with its left, right, and top
previews. Formal 2048 processing then reuses the selected source exactly.

This case is not merged with the procedural formal A/B/C pool. Its purpose is to
test the complete human-Agent interaction, real image-generation backend,
checkpoint persistence, candidate-local retry, and final export contract.

\subsection{Adapter transfer case}

The M4A4 case reuses the same theme assets and core algorithms with zero additional
image-generation calls. Only the adapter, official mesh hash, canonical axes, UV
topology, and seam calibration change. The M4A4 mesh contains 45,422 triangles,
39,881 UV coordinates, 1,747 UV islands, and 11,341 seam pairs. The test asks
whether the pipeline executes and where the AK-47 thresholds cease to transfer.

\subsection{Evaluation criteria}

The report uses the following primary measures:

\begin{itemize}
  \item asset-seam error $\assetseam$ (lower is better);
  \item multi-view score $\multiview$ (higher is better);
  \item multi-view retention relative to the uncorrected candidate;
  \item locked Agent score for Route C accept/rollback decisions;
  \item changed-pixel fraction and changed pixels outside the target halo;
  \item successful completion of the three blocking human checkpoints;
  \item final texture resolution, colour mode, export identity, and file hash.
\end{itemize}

Automatic readability is reported only for explanation. It is not a substitute
for a human preference study, and no claim is made that it measures general
aesthetic quality.

\subsection{Implementation and reproducibility}

The implementation uses Python 3.13, NumPy, Pillow, OpenCV, SciPy, PyTorch,
Transformers, Diffusers, and Streamlit. The local diffusion acceptance used an
NVIDIA RTX 4060 Laptop GPU with 8 GB VRAM. The repository contains a deterministic
procedural backend, a pinned SD-Turbo backend, provider adapters, the Agent
runtime, command-line scripts, and a thin Streamlit client. At the report freeze,
100 unit tests, Python byte-code compilation, configuration parsing, and headless
Streamlit startup had passed.

The public repository is \url{\repositoryurl}. Valve geometry, model weights, and
API keys are intentionally excluded. Setup documentation explains how authorized
users obtain external assets and configure local credentials.
